\documentclass[12pt,a4paper]{article}

\usepackage[utf8]{inputenc}
\usepackage[T1]{fontenc}
\usepackage{lmodern}
\usepackage[portuguese]{babel}

\usepackage{listings}
\usepackage{xcolor}
\usepackage{amsmath}
\usepackage{amssymb}
\usepackage{amsfonts}
\usepackage{graphicx}

\lstset{
    basicstyle=\ttfamily\small,
    keywordstyle=\color{blue}\bfseries, 
    commentstyle=\color{green!70!black},
    stringstyle=\color{red}, 
    numbers=left, 
    numberstyle=\tiny\color{gray}, 
    frame=single, 
    breaklines=true, 
    showstringspaces=false 
}

\usepackage[margin=2.5cm]{geometry}

\begin{document}

\begin{center}
    {\LARGE MC521 - Relatório X} \\ [2,5em]
    {\large Júlio da Silva Telles} \\ [1,0em]
    {\large 22 de Junho de 2026} \\ [1,5em]
\end{center}

\vspace{1cm} 


\section{Dijkstra? (CodeForces - 20C)}
O problema pede explicitamente para encontrar o caminho mais curto entre o vértice $1$ e o vértice $N$ em um grafo ponderado e não-direcionado. O diferencial prático é que não basta apenas calcular a distância final: é necessário recuperar e imprimir a sequência exata de vértices que compõem esse caminho.

\subsection{Solução}
A solução baseia-se no Algoritmo de Dijkstra aliado a um vetor de rastreamento de predecessores para permitir o \textit{backtracking} da rota. A complexidade de tempo atinge $\mathcal{O}(M \log N)$. O fluxo funciona da seguinte maneira:

\begin{enumerate}
    \item \textbf{Exploração do Grafo:} Utilizamos o Dijkstra padrão auxiliado por uma fila de prioridades (\texttt{priority\_queue}) orientada a valores mínimos (\texttt{greater<ii>}), extraindo sempre o nó com a menor distância acumulada.
    \item \textbf{Rastreamento (Predecessores):} Criamos um vetor \texttt{pred} inicializado com $-1$. Sempre que encontramos um caminho mais curto para um nó vizinho (\texttt{v.first}) passando pelo nó atual (\texttt{u}), atualizamos a distância mínima e salvamos a informação de que "viemos de $u$" fazendo \texttt{pred[v.first] = u}.
    \item \textbf{Reconstrução do Caminho:} Após o fim do algoritmo, checamos se o vértice $N$ foi alcançado (\texttt{dist[n] != INF}). Se sim, utilizamos a função recursiva \texttt{print\_path} para percorrer o vetor \texttt{pred} de trás para frente (começando em $N$ e indo até o $1$). A recursão garante que os nós sejam impressos na ordem correta da origem ao destino. Se $N$ for inalcançável, imprime-se $-1$.
\end{enumerate}

\newpage

\section{Road Construction (Aizu - 2249)}
O problema descreve a necessidade de reconstruir as estradas de um país partindo da capital (cidade 1) para todas as outras cidades, mantendo a distância mínima original entre a capital e cada cidade. Como cada estrada possui um custo de construção, o objetivo é selecionar um subconjunto de estradas que satisfaça a condição das distâncias mínimas e, ao mesmo tempo, minimize o custo total da obra.

\subsection{Solução}
Este problema exige uma variação do Algoritmo de Dijkstra para encontrar a Árvore de Caminhos Mínimos (\textit{Shortest Path Tree}), ponderando pelo custo em caso de empates. A complexidade de tempo é de $\mathcal{O}(M \log N)$, onde $M$ é o número de estradas e $N$ o número de cidades. O algoritmo segue os passos:

\begin{enumerate}
    \item \textbf{Cálculo das Distâncias Mínimas:} Primeiro, executamos o algoritmo de Dijkstra tradicional partindo da capital (nó 1) utilizando exclusivamente as \textit{distâncias} como peso das arestas. Isso preenche o vetor \texttt{dist} com a distância mínima exata da capital para cada cidade do país.
    \item \textbf{Otimização de Custos (Tie-breaking):} Após conhecer a distância ótima para cada cidade, precisamos escolher qual estrada construir para conectá-la à malha. Iteramos sobre cada cidade $i$ (de $2$ até $N$) e analisamos todas as suas estradas adjacentes.
    \item \textbf{Seleção da Aresta:} Para cada vizinho $v$, verificamos se a estrada pertence a um caminho mínimo válido através da equação \texttt{dist[v] + peso == dist[i]}. Caso seja uma estrada válida, atualizamos o menor custo de construção necessário para conectar a cidade $i$ (\texttt{min\_cost}). Por fim, somamos esse menor custo à variável \texttt{total}.
\end{enumerate}

\end{document}